\section{Machine learning architectures}
\label{sec:architectures}

\subsection{System architecture}

Figure~\ref{fig:pipeline} shows the path of one request. The application is
split into three processes that share only files and HTTP: the notebooks
produce the model bundles, the API serves them, and the desktop application
consumes the API. Inside the application every stage is behind a protocol
(an abstract interface), so the Whisper backend, the Piper backend, the ONNX
emotion detector, the Gemini parser and the HTTP client are each
replaceable without touching the orchestrator.

\begin{figure}[htbp]
  \centering
  \begin{tikzpicture}[
      x=3.1cm, y=1.9cm,
      every node/.style={font=\small},
      box/.style={draw, rounded corners, align=center, minimum height=1.15cm, text width=2.7cm},
      ml/.style={box, fill=blue!8},
      io/.style={box, fill=gray!10},
      arr/.style={-{Latex[length=2mm]}, thick},
    ]
    \node[io] (mic) at (0,2) {Microphone\\ \scriptsize sounddevice, 16 kHz};
    \node[ml] (stt) at (1,2) {Speech-to-text\\ \scriptsize Whisper \emph{small}, int8};
    \node[box] (wake) at (2,2) {Wake word\\ \scriptsize \emph{JarvisTEC} + mishearings};
    \node[ml] (intent) at (3,2) {Intent parser\\ \scriptsize Gemini $\to$ rules};
    \node[box] (api) at (3,1) {REST API\\ \scriptsize FastAPI, \code{/predict}};
    \node[ml] (models) at (4,1) {10 models\\ \scriptsize joblib bundles};
    \node[box] (reply) at (2,1) {Reply builder\\ \scriptsize template per model};
    \node[box] (tone) at (1,1) {Tone adapter\\ \scriptsize prefix/suffix};
    \node[ml] (tts) at (0,1) {Text-to-speech\\ \scriptsize Piper es\_MX};
    \node[io] (cam) at (0,0) {Webcam\\ \scriptsize OpenCV, $\sim$15 fps};
    \node[box] (face) at (1,0) {Face locator\\ \scriptsize Haar cascade};
    \node[ml] (emo) at (2,0) {Emotion CNN\\ \scriptsize ONNX, 48$\times$48, 7 classes};
    \node[box] (avatar) at (3,0) {Avatar\\ \scriptsize expression per emotion};

    \draw[arr] (mic) -- (stt);
    \draw[arr] (stt) -- (wake);
    \draw[arr] (wake) -- (intent);
    \draw[arr] (intent) -- (api);
    \draw[arr, <->] (api) -- (models);
    \draw[arr] (api) -- (reply);
    \draw[arr] (reply) -- (tone);
    \draw[arr] (tone) -- (tts);
    \draw[arr] (cam) -- (face);
    \draw[arr] (face) -- (emo);
    \draw[arr] (emo) -- (avatar);
    \draw[arr] (emo.north) -- ++(0,0.22) -| (tone.south);
  \end{tikzpicture}
  \caption{One request through \jarvis{}. Shaded boxes are learned models. Typed text enters at the wake-word stage and follows the same path.}
  \label{fig:pipeline}
\end{figure}

\subsection{Perception agents}

\paragraph{Facial emotion agent.}
Every 2 seconds the camera worker takes the current frame, converts it to
grayscale and runs the Viola--Jones Haar cascade
(\code{haarcascade\_frontalface\_alt}) to locate the largest frontal face,
rejecting boxes under $60\times60$ pixels to avoid glare and texture
patches. The crop is resized to $48\times48$, scaled to $[0,1]$ and fed to a
VGG-style convolutional network: five blocks of
$3\times3$ convolutions with ReLU (64, 128, 256, 512 and 512 filters, two or
three convolutions per block), each block closed by batch normalisation and
$2\times2$ max pooling, then a dense layer of 128 units with dropout 0.2 and
a softmax over the seven classes. The network was trained on FER-2013 and reused from an earlier course project; for this
project the Keras weights were exported once to ONNX so that inference runs
with \code{onnxruntime} on the CPU with two threads and does not starve the
interface. The winning class always drives the avatar and the tone; with
seven classes a webcam reading rarely exceeds 50\,\% probability, so the
confidence is shown for information only. The reading is repeated every two
seconds and the face box is redrawn on every frame in between.

\paragraph{Speech agent.}
\code{Ctrl+Space} or the microphone button opens a 16\,kHz mono recording;
pressing again closes it. The audio goes to Whisper \emph{small} through \code{faster-whisper} (CTranslate2), loaded
in int8 on the CPU with Spanish fixed as the language. The transcript then
follows exactly the same path as typed text: wake-word check, intent
parsing, prediction, reply. Whisper's typical mishearings of the wake word
are part of the wake-word table, and the Gemini prompt is told the text may
come from speech so it can recover \emph{días} from \emph{dientes}.

\paragraph{Intent parser.}
The parser turns a Spanish sentence into
\code{(model\_name, fields, missing)}. With a key configured, Gemini is asked
for a JSON object under a system instruction that carries the catalogue of
the ten models (a one-line description, Spanish synonyms and per-model
inference rules) and the JSON schema of every request body, which the
application reads from the API's OpenAPI document at startup. Models are
tried in a chain and a model whose quota is exhausted (429) or that was
retired (404) is skipped for the rest of the session. Without a key, or on
any Gemini failure, a rule-based parser with one field extractor per model
(keyword lookup tables, Spanish number words, dates, airline and city codes)
produces the same structure. Whichever parser answered is shown in the
status strip.

\paragraph{Dialogue.}
The orchestrator records the turn, resolves the intent against the pending
one (so an answer to ``which airline?'' completes the flight request from
the previous turn), asks for the next missing required field if any, or
posts the request body to the API and renders the answer through a per-model
reply template. The tone adapter wraps the text according to the last
emotion, and the speaking worker synthesises it with Piper
(\code{es\_MX-claude-high}, 22\,kHz). Starting a new recording interrupts the
voice.

\subsection{The ten models}

Table~\ref{tab:models} lists the ten problems, their datasets, the technique
required by the course for each and the target. Each one is developed in
its own notebook with the six CRISP-DM phases: problem
analysis, data understanding, exploration and preparation, modelling,
evaluation and deployment (the saved bundle plus a reload test). Every
notebook builds the literal model the task asks for first, and when the data
call for a deviation (a logarithm, a regrouped target, a balanced class
weight) that second model is built next to it and the choice is defended
with the metrics of both.

\begin{table}[htbp]
  \centering
  \small
  \begin{tabularx}{\textwidth}{@{}c>{\ttfamily}lLLL@{}}
    \toprule
    \# & Notebook & Dataset (source) & Technique & Target \\
    \midrule
    1 & bitcoin\_price\_regression & Bitcoin daily prices 2013--2017 (Kaggle, team-ai) & ARIMA & closing price, next days \\
    2 & car\_price\_regression & 301 used cars, India (Kharwal) & multiple linear regression & selling price (lakhs) \\
    3 & telecom\_churn\_classification & IBM Telco churn, 7\,043 customers & neural network (MLP) & churn yes/no \\
    4 & wine\_quality\_classification & Vinho verde, 6\,497 wines & $k$ nearest neighbours & quality band \\
    5 & flight\_delay\_regression & BTS airline delays 2008, 1.9\,M flights & decision tree & arrival delay (min) \\
    6 & stroke\_classification & Stroke prediction, 5\,110 people (fedesoriano) & logistic regression & stroke yes/no \\
    7 & cirrhosis\_type\_classification & Mayo PBC trial, 418 patients & support vector machine & cirrhosis stage \\
    8 & avocado\_price\_regression & Hass Avocado Board retail scans 2015--2018 & random forest & average price (USD) \\
    9 & hepatitis\_type\_classification & HCV blood panels, 615 people & Gaussian naive Bayes & donor / hepatitis / fibrosis / cirrhosis \\
    10 & body\_fat\_regression & 252 men & multiple linear regression & body fat (\%) \\
    \bottomrule
  \end{tabularx}
  \caption{The ten models, in the order of the repository. All datasets are downloaded by the notebooks through \code{kagglehub} or a public URL.}
  \label{tab:models}
\end{table}

\subsection{Serving architecture}

The API is a FastAPI service with one adapter per model. Each adapter knows
how to turn a validated request body into the row the estimator expects
(filling optional fields with the notebook's fill value, applying the
notebook's encodings and transformations) and how to turn the estimator's
output into the response (label, probability and threshold for classifiers;
value and units for regressions; a list of \code{\{date, price, lower,
upper\}} for the forecast). Requests are validated by Pydantic schemas that
mirror the notebook's predictor list, and regressions attach a warning when
an input falls outside the training range, since a fitted line has no
information beyond it. Endpoints, request bodies and example calls are listed
in Appendix~\ref{app:api}.
